\documentclass[instructions.tex]{subfiles}
\begin{document}
 
\chapter{Pour faire une bonne confession , il faut avoir le propos de changer de vie.}
 
\primo Un arbre qui porte des feuilles et des fleurs est un bel arbre; mais s'il ne porte point de fruit, il n'est pas un bon arbre, il mérite le feu\footnote{Omnis arbor quæ non facit fructum bonum , excidetur et in ignem mittetur. \textit{Matth.} 3.}. En vous confessant, vous faites des résolutions, de belles promesses, ce sont là les feuilles et les fleurs de l'arbre; mais s'il n'y a point de fruits de pénitence et d'amendement, vous êtes un arbre mort. Vous les connoîtrez à leurs fruits, dit Jésus-Christ.
 
Par la confession, vous professez votre christianisme et votre foi; mais de quoi sert cette profession de foi, si vous ne changez de vie? Ne savez-vous pas , dit saint Jacques , que la foi sans les œuvres est morte , de même que le corps sans l'âme est un corps mort? Que devez-vous donc penser de vos confessions?
 
\secundo On se confesse pour recevoir le pardon de ses péchés par l'absolution d'un prêtre revêtu de l'autorité de Jésus-Christ; mais vous aveugleriez-vous jusqu'à vous persuader que Dieu pardonne des péchés qu'on ne veut pas quitter? Or, veut-on quitter ses péchés quand on veut toujours les commettre? et ne veut-on pas toujours les commettre, quand on les commet toujours en effet? Rendez-vous ici justice. Quel péché avez-vous quitté? quelle occasion avez-vous évitée? Avez-vous réparé vos scandales , vos médisances, vos injustices et vos larcins? avez-vous réglé votre conduite et votre famille? Quel changement voit-on en vous?
 
Vous disiez que vous vouliez changer, et vous l'avez promis; mais votre volonté a-t-elle été sincère? Une promesse est toujours suspecte, lorsqu'elle est sans effet. Dire qu'on veut changer, et vivre toujours de même, c'est dire en spéculation qu'on veut se sauver, et vouloir se damner dans la pratique; ou plutôt, c'est se moquer de Dieu. Celui-là , dit saint Augustin, est un moqueur et un faux pénitent, qui continue de faire ce qu'il avait détesté et ce qu'il avait promis de ne plus faire\footnote{Ille irrisor est, non pœnitens, qui adhuc agit quod pœnituit. \textit{S. Aug.}}. Combien de gens seront étonnés à la mort, lorsque Dieu leur reprochera tant de confessions faites sans propos d'amendement!
 
\section{Résolutions}
 
\primo Enfin je ne promettrai plus pour ne pas tenir mes promesses.
 
\secundo Je corrigerai donc dès à présent ce que j'ai déjà tant dit de fois que je voulois corriger. 

\prayer{Pardonnez-moi, ô mon Dieu! les mauvaises confessions que j'ai faites par défaut de bon propos; aidez-moi à les réparer, de concert avec celui que vous avez chargé de la conduite de ma conscience.}
 
\end{document}
